% Generado por revision/generar_e1.py 3 en orden de cita; no editar a mano.
\begin{thebibliography}{44}
\bibitem{Mockus2000} A. Mockus and D. M. Weiss, ``Predicting risk of software changes,'' Bell Labs Tech. J., vol. 5, no. 2, pp. 169--180, Apr.--Jun. 2000, doi: \mbox{10.1002/bltj.2229}.
\bibitem{Kamei2013} Y. Kamei, E. Shihab, B. Adams, A. E. Hassan, A. Mockus, A. Sinha, and N. Ubayashi, ``A large-scale empirical study of just-in-time quality assurance,'' IEEE Trans. Softw. Eng., vol. 39, no. 6, pp. 757--773, Jun. 2013, doi: \mbox{10.1109/TSE.2012.70}.
\bibitem{Sliwerski2005} J. Śliwerski, T. Zimmermann, and A. Zeller, ``When do changes induce fixes?'' in Proc. Int. Workshop Mining Softw. Repositories (MSR '05), St. Louis, MO, USA, 2005, pp. 1--5, doi: \mbox{10.1145/1083142.1083147}.
\bibitem{Lyu2024} Y. Lyu, H. J. Kang, R. Widyasari, J. Lawall, and D. Lo, ``Evaluating SZZ implementations: An empirical study on the Linux kernel,'' IEEE Trans. Softw. Eng., vol. 50, no. 9, pp. 2219--2239, Sep. 2024, doi: \mbox{10.1109/TSE.2024.3406718}.
\bibitem{KernelAdd} S. Levin, ``docs: add AI Coding Assistants documentation,'' Linux kernel, commit 78d979db6cef, Jan. 6, 2026. [Online]. Available: \url{https://github.com/torvalds/linux/commit/78d979db6cef}
\bibitem{KernelSimpl} C. Brauner, ``coding-assistants: simplify attribution,'' Linux kernel, commit 816d9992d9ed, Aug. 3, 2026. [Online]. Available: \url{https://github.com/torvalds/linux/commit/816d9992d9ed}
\bibitem{KernelRemove} S. Levin, ``docs: remove Assisted-by references,'' Linux kernel, commit 1e035297ffda, Oct. 14, 2026. [Online]. Available: \url{https://github.com/torvalds/linux/commit/1e035297ffda}
\bibitem{Studer2021} S. Studer, T. B. Bui, C. Drescher, A. Hanuschkin, L. Winkler, S. Peters, and K.-R. M\"uller, ``Towards CRISP-ML(Q): A machine learning process model with quality assurance methodology,'' Mach. Learn. Knowl. Extr., vol. 3, no. 2, pp. 392--413, 2021, doi: 10.3390/make3020020.
\bibitem{MaesBermejo2024} M. Maes Bermejo, J. M. González-Barahona, M. Gallego, and G. Robles, ``A dataset of Linux kernel commits'' (v3), Zenodo, 2024, doi: 10.5281/zenodo.10654193.
\bibitem{Page2021} M. J. Page \emph{et al.}, ``The PRISMA 2020 statement: An updated guideline for reporting systematic reviews,'' BMJ, vol. 372, Art. no. n71, 2021, doi: 10.1136/bmj.n71.
\bibitem{Zeng2021} Z. Zeng, Y. Zhang, H. Zhang, and L. Zhang, ``Deep just-in-time defect prediction: How far are we?'' in Proc. 30th ACM SIGSOFT Int. Symp. Softw. Testing Anal. (ISSTA), Virtual, Denmark, 2021, pp. 427--438, doi: \mbox{10.1145/3460319.3464819}.
\bibitem{Ni2022} C. Ni, W. Wang, K. Yang, X. Xia, K. Liu, and D. Lo, ``The best of both worlds: Integrating semantic features with expert features for defect prediction and localization,'' in Proc. 30th ACM Joint Eur. Softw. Eng. Conf. Symp. Found. Softw. Eng. (ESEC/FSE), Singapore, 2022, pp. 672--683, doi: \mbox{10.1145/3540250.3549165}.
\bibitem{Liu2026} Y. Liu, K. Jabre, R. Abreu, Z. J. Carmichael, V. Murali, A. Patel, J. Ge, W. Sun, C. Zhang, A. Mockus, D. Khavari, P. C. Rigby, and N. Nagappan, ``A preliminary study on explaining risk of code changes using LLM-based prediction models,'' in Proc. 3rd ACM Int. Conf. AI-Powered Softw. (AIware '26), Montreal, QC, Canada, 2026, pp. 259--267, doi: \mbox{10.1145/3805760.3814916}.
\bibitem{JITLine2021} C. Pornprasit and C. K. Tantithamthavorn, ``JITLine: A simpler, better, faster, finer-grained just-in-time defect prediction,'' in Proc. IEEE/ACM 18th Int. Conf. Mining Softw. Repositories (MSR), 2021, pp. 369--379, doi: \mbox{10.1109/MSR52588.2021.00049}.
\bibitem{PyExplainer2021} C. Pornprasit, C. Tantithamthavorn, J. Jiarpakdee, M. Fu, and P. Thongtanunam, ``PyExplainer: Explaining the predictions of just-in-time defect models,'' in Proc. 36th IEEE/ACM Int. Conf. Automated Softw. Eng. (ASE), 2021, pp. 407--418, doi: \mbox{10.1109/ASE51524.2021.9678763}.
\bibitem{Aftab2023} S. Aftab, S. Abbas, T. M. Ghazal, M. Ahmad, H. A. Hamadi, C. Y. Yeun, and M. A. Khan, ``A cloud-based software defect prediction system using data and decision-level machine learning fusion,'' Mathematics, vol. 11, no. 3, Art. no. 632, 2023, doi: \mbox{10.3390/math11030632}.
\bibitem{Dey2025} R. Dey, J. Piri, B. Acharya, P. P. Das, V. C. Gerogiannis, and A. Kanavos, ``A hybrid evolutionary fuzzy ensemble approach for accurate software defect prediction,'' Mathematics, vol. 13, no. 7, Art. no. 1140, 2025, doi: \mbox{10.3390/math13071140}.
\bibitem{McIntosh2018} S. McIntosh and Y. Kamei, ``Are fix-inducing changes a moving target? A longitudinal case study of just-in-time defect prediction,'' IEEE Trans. Softw. Eng., vol. 44, no. 5, pp. 412--428, May 2018, doi: \mbox{10.1109/TSE.2017.2693980}.
\bibitem{Cabral2019} G. G. Cabral, L. L. Minku, E. Shihab, and S. Mujahid, ``Class imbalance evolution and verification latency in just-in-time software defect prediction,'' in Proc. IEEE/ACM 41st Int. Conf. Softw. Eng. (ICSE), Montreal, QC, Canada, 2019, pp. 666--676, doi: \mbox{10.1109/ICSE.2019.00076}.
\bibitem{Olewicki2024} D. Olewicki, S. Habchi, M. Nayrolles, M. Faramarzi, S. Chandar, and B. Adams, ``On the costs and benefits of adopting lifelong learning for software analytics -- Empirical study on brown build and risk prediction,'' in Proc. 46th Int. Conf. Softw. Eng.: Softw. Eng. Pract. (ICSE-SEIP), Lisbon, Portugal, 2024, pp. 275--286, doi: \mbox{10.1145/3639477.3639717}.
\bibitem{Nguyen2025} D. Nguyen, T. Le-Cong, T. H. M. Le, M. A. Babar, and Q.-T. Huynh, ``Toward realistic evaluations of just-in-time vulnerability prediction,'' in Proc. IEEE Int. Conf. Softw. Maint. Evol. (ICSME), Auckland, New Zealand, 2025, pp. 1--10, doi: \mbox{10.1109/ICSME64153.2025.00023}.
\bibitem{Mazloomzadeh2026} I. Mazloomzadeh, M. M. Morovati, and F. Khomh, ``How do AI coding agents contribute to software development? An empirical study of agentic pull requests,'' arXiv:2607.21832, Jul. 2026, doi: \mbox{10.48550/arXiv.2607.21832}.
\bibitem{Herbold2022} S. Herbold, A. Trautsch, F. Trautsch, and B. Ledel, ``Problems with SZZ and features: An empirical study of the state of practice of defect prediction data collection,'' Empir. Softw. Eng., vol. 27, no. 2, Art. no. 42, 2022, doi: \mbox{10.1007/s10664-021-10092-4}.
\bibitem{Falessi2022} D. Falessi, A. Ahluwalia, and M. Di Penta, ``The impact of dormant defects on defect prediction: A study of 19 Apache projects,'' ACM Trans. Softw. Eng. Methodol., vol. 31, no. 1, pp. 1--26, Jan. 2022, doi: \mbox{10.1145/3467895}.
\bibitem{Jiarpakdee2021} J. Jiarpakdee, C. K. Tantithamthavorn, and J. Grundy, ``Practitioners' perceptions of the goals and visual explanations of defect prediction models,'' in Proc. IEEE/ACM 18th Int. Conf. Mining Softw. Repositories (MSR), 2021, pp. 432--443, doi: \mbox{10.1109/MSR52588.2021.00055}.
\bibitem{Sahar2024} H. Sahar, A. A. Bangash, A. Hindle, and D. Barbosa, ``IRJIT: A simple, online, information retrieval approach for just-in-time software defect prediction,'' Empir. Softw. Eng., vol. 29, no. 5, Art. no. 131, 2024, doi: \mbox{10.1007/s10664-024-10514-z}.
\bibitem{Yu2024FoX} J. Yu, M. Fu, A. Ignatiev, C. Tantithamthavorn, and P. Stuckey, ``A formal explainer for just-in-time defect predictions,'' ACM Trans. Softw. Eng. Methodol., vol. 33, no. 7, pp. 1--31, Sep. 2024, doi: \mbox{10.1145/3664809}.
\bibitem{Chen2018MULTI} X. Chen, Y. Zhao, Q. Wang, and Z. Yuan, ``MULTI: Multi-objective effort-aware just-in-time software defect prediction,'' Inf. Softw. Technol., vol. 93, pp. 1--13, Jan. 2018, doi: \mbox{10.1016/j.infsof.2017.08.004}.
\bibitem{Kang2021} J. Kang, S. Kwon, D. Ryu, and J. Baik, ``HASPO: Harmony search-based parameter optimization for just-in-time software defect prediction in maritime software,'' Appl. Sci., vol. 11, no. 5, Art. no. 2002, 2021, doi: \mbox{10.3390/app11052002}.
\bibitem{Mostafa2025} S. Mostafa, S. T. Cynthia, B. Roy, and D. Mondal, ``Feature transformation for improved software bug detection and commit classification,'' J. Syst. Softw., vol. 219, Art. no. 112205, Jan. 2025, doi: \mbox{10.1016/j.jss.2024.112205}.
\bibitem{Chen2025} F. Chen, Y. Ke, X. Liu, and Q. Wei, ``Multi-project just-in-time software defect prediction based on multi-task learning for mobile applications,'' in Proc. IEEE Conf. Softw. Testing, Verification Validation (ICST), 2025, pp. 127--137, doi: \mbox{10.1109/ICST62969.2025.10989004}.
\bibitem{Shehab2023} M. A. Shehab, A. Hamou-Lhadj, and V. S. P. Gunda, ``JITBoost: Boosting just-in-time defect prediction using Boolean combination of classifiers,'' in Proc. IEEE 23rd Int. Conf. Softw. Qual., Rel. Secur. (QRS), 2023, pp. 95--104, doi: \mbox{10.1109/QRS60937.2023.00019}.
\bibitem{Chitsazian2026} Z. Chitsazian, S. Sedighian Kashi, and A. Nikanjam, ``Detecting concept drift in just-in-time software defect prediction using model interpretation,'' Int. J. Syst. Assur. Eng. Manag., vol. 17, no. 9, pp. 2701--2722, 2026, doi: \mbox{10.1007/s13198-025-03122-7}.
\bibitem{Patel2024} H. J. Patel, B. Adams, and A. E. Hassan, ``Post deployment recycling of machine learning models,'' Empir. Softw. Eng., vol. 29, no. 4, Art. no. 100, 2024, doi: \mbox{10.1007/s10664-024-10492-2}.
\bibitem{Tian2026} D. Tian, L. Song, and X. Su, ``Just-in-time quality assurance with assistance of automated defect prediction and code testing,'' ACM Trans. Softw. Eng. Methodol., vol. 35, no. 9, pp. 1--46, 2026, doi: \mbox{10.1145/3779653}.
\bibitem{Song2023} L. Song and L. L. Minku, ``A procedure to continuously evaluate predictive performance of just-in-time software defect prediction models during software development,'' IEEE Trans. Softw. Eng., vol. 49, no. 2, pp. 646--666, Feb. 2023, doi: \mbox{10.1109/TSE.2022.3158831}.
\bibitem{Veras2025} R. C. Veras, G. G. Cabral, and A. L. I. de Oliveira, ``Semantic SZZ: Mitigating the impact of misclassified corrective changes in just-in-time software defect prediction,'' in Proc. IEEE Int. Conf. Syst., Man, Cybern. (SMC), 2025, pp. 769--774, doi: \mbox{10.1109/SMC58881.2025.11342537}.
\bibitem{Friedman2001} J. H. Friedman, ``Greedy function approximation: A gradient boosting machine,'' Ann. Statist., vol. 29, no. 5, pp. 1189--1232, Oct. 2001, doi: \mbox{10.1214/aos/1013203451}.
\bibitem{Ruspini1969} E. H. Ruspini, ``A new approach to clustering,'' Inf. Control, vol. 15, no. 1, pp. 22--32, Jul. 1969, doi: \mbox{10.1016/S0019-9958(69)90591-9}.
\bibitem{Herrera2008} F. Herrera, ``Genetic fuzzy systems: taxonomy, current research trends and prospects,'' Evol. Intell., vol. 1, no. 1, pp. 27--46, Mar. 2008, doi: \mbox{10.1007/s12065-007-0001-5}.
\bibitem{Ishibuchi2001} H. Ishibuchi, T. Nakashima, and T. Murata, ``Three-objective genetics-based machine learning for linguistic rule extraction,'' Inf. Sci., vol. 136, no. 1--4, pp. 109--133, Aug. 2001, doi: \mbox{10.1016/S0020-0255(01)00144-X}.
\bibitem{Saito2015} T. Saito and M. Rehmsmeier, ``The precision-recall plot is more informative than the ROC plot when evaluating binary classifiers on imbalanced datasets,'' PLOS ONE, vol. 10, no. 3, Art. no. e0118432, Mar. 2015, doi: \mbox{10.1371/journal.pone.0118432}.
\bibitem{Dietterich1998} T. G. Dietterich, ``Approximate statistical tests for comparing supervised classification learning algorithms,'' Neural Comput., vol. 10, no. 7, pp. 1895--1923, Oct. 1998, doi: \mbox{10.1162/089976698300017197}.
\bibitem{Demsar2006} J. Demšar, ``Statistical comparisons of classifiers over multiple data sets,'' J. Mach. Learn. Res., vol. 7, pp. 1--30, 2006. [Online]. Available: \url{https://jmlr.org/papers/v7/demsar06a.html}
\end{thebibliography}
